\chapter{Arc Length of the Planar Archimedean Spiral: Two Methods}
\label{sec:bogenlaenge_spirale}

In this chapter the arc length of the Archimedean spiral
$r(\theta) = \frac{6}{\pi^2}\,\theta$ is computed by two entirely different
methods, and the results are compared numerically. This serves both the
mutual verification and the conceptual clarification of how the
$\theta$-parametrisation and the $p$-parametrisation are related.
Throughout the numerical tables, the decimal point is used consistently.

The intervals considered are the minor steps between the
integer indices $n \equiv 1, 5 \pmod 6$:
\[
(n_1, n_2) \in
\{(5,7),\;(11,13),\;(17,19),\;(23,25),\;(29,31),\;(35,37),\;(41,43)\}.
\]
In each case $n_1 \equiv 5 \pmod 6$ and $n_2 \equiv 1 \pmod 6$, so these
are always minor steps with $\Delta n = 2$.

% ─────────────────────────────────────────────────────────────────────────────
\section{Method A: Direct Computation via the Closed-Form Arc Length Formula}
\label{sec:methode_A}
% ─────────────────────────────────────────────────────────────────────────────

\subsection{Derivation of the Formula}

For a curve in polar coordinates $r = r(\theta)$, the
arc length element is
\[
ds = \sqrt{r(\theta)^2 + \left(\frac{dr}{d\theta}\right)^2}\,d\theta.
\]
With $r(\theta) = a\theta$ and $\frac{dr}{d\theta} = a$ (where $a = \frac{6}{\pi^2}$):
\[
ds = \sqrt{a^2\theta^2 + a^2}\,d\theta = a\sqrt{\theta^2 + 1}\,d\theta.
\]

The arc length integral between two angles $\theta_1$ and $\theta_2$ is thus:
\[
L_A(\theta_1, \theta_2)
= a\int_{\theta_1}^{\theta_2}\sqrt{\theta^2 + 1}\,d\theta.
\]

This integral admits a closed form. Using the standard formula
\[
\int \sqrt{u^2+1}\,du
= \frac{1}{2}\Bigl[u\sqrt{u^2+1} + \operatorname{arcsinh}(u)\Bigr] + C,
\]
the antiderivative is $F(\theta) = \frac{a}{2}\bigl[\theta\sqrt{\theta^2+1}
+ \operatorname{arcsinh}(\theta)\bigr]$, yielding the exact, closed-form arc length:

\begin{satz}[Arc length of the Archimedean spiral, closed form]
\label{satz:bogenlaenge_geschlossen}
For $r(\theta) = a\theta$ with $a = \frac{6}{\pi^2}$:
\begin{equation}
\boxed{
L_A(\theta_1, \theta_2)
= \frac{a}{2}\Bigl[\theta\sqrt{\theta^2+1} + \operatorname{arcsinh}(\theta)\Bigr]_{\theta_1}^{\theta_2}.
}
\end{equation}
\end{satz}

\subsection{Assignment of Angular Values to the Indices}

The angular values $\theta_n$ at the integer points of the helix are
\[
\theta_n = \frac{\pi}{3}\,n + \frac{\pi}{4} = \frac{(4n+3)\pi}{12}.
\]
As exact fractions of $\pi$:

\begin{center}
\begin{tabular}{ccccc}
\toprule
$n_1$ & $\theta_{n_1} = \frac{(4n_1+3)\pi}{12}$ & $n_2$ & $\theta_{n_2} = \frac{(4n_2+3)\pi}{12}$ & $\Delta\theta = \frac{2\pi}{3}$ \\
\midrule
 5 & $\frac{23\pi}{12}$ &  7 & $\frac{31\pi}{12}$ & $\frac{8\pi}{12} = \frac{2\pi}{3}$ \\
11 & $\frac{47\pi}{12}$ & 13 & $\frac{55\pi}{12}$ & $\frac{2\pi}{3}$ \\
17 & $\frac{71\pi}{12}$ & 19 & $\frac{79\pi}{12}$ & $\frac{2\pi}{3}$ \\
23 & $\frac{95\pi}{12}$ & 25 & $\frac{103\pi}{12}$ & $\frac{2\pi}{3}$ \\
29 & $\frac{119\pi}{12}$ & 31 & $\frac{127\pi}{12}$ & $\frac{2\pi}{3}$ \\
35 & $\frac{143\pi}{12}$ & 37 & $\frac{151\pi}{12}$ & $\frac{2\pi}{3}$ \\
41 & $\frac{167\pi}{12}$ & 43 & $\frac{175\pi}{12}$ & $\frac{2\pi}{3}$ \\
\bottomrule
\end{tabular}
\end{center}

\begin{bemerkung}[Constant angular spacing]
All minor steps (i.e.\ $\Delta n = 2$) produce the same angular spacing
$\Delta\theta = \frac{2\pi}{3}$, since
\[
  \theta_{n+2} - \theta_n
  = \frac{\pi}{3}(n+2) + \frac{\pi}{4} - \frac{\pi}{3}n - \frac{\pi}{4}
  = \frac{2\pi}{3}.
\]
The spiral traverses exactly $\frac{1}{3}$ of a full revolution per minor step.
The angular spacing is thus identical for all intervals under consideration;
the arc length nevertheless grows monotonically, since for equal $\Delta\theta$
the arc on a larger radius is longer.
\end{bemerkung}

\subsection{Numerical Results for Method A}

\begin{table}[H]
\centering
\caption{Arc length Method~A: closed-form formula
  $L_A = \frac{a}{2}[\theta\sqrt{\theta^2+1}+\operatorname{arcsinh}(\theta)]_{\theta_1}^{\theta_2}$
  with $a = 6/\pi^2$.}
\renewcommand{\arraystretch}{1.2}
\begin{tabular}{ccccc}
\toprule
$(n_1, n_2)$ & $\theta_1$ (num.) & $\theta_2$ (num.) & $L_A$ \\
\midrule
$( 5,  7)$ & 6.02138592 &  8.11578102 &  9.09026487 \\
$(11, 13)$ & 12.30457123 & 14.39896633 & 17.04771118 \\
$(17, 19)$ & 18.58775653 & 20.68215164 & 25.03243246 \\
$(23, 25)$ & 24.87094184 & 26.96533694 & 33.02456693 \\
$(29, 31)$ & 31.15412715 & 33.24852225 & 41.01977219 \\
$(35, 37)$ & 37.43731245 & 39.53170755 & 49.01654352 \\
$(41, 43)$ & 43.72049776 & 45.81489286 & 57.01422134 \\
\bottomrule
\end{tabular}
\end{table}

\begin{bemerkung}[Growth rate of the arc length]
The arc lengths grow approximately proportionally to the mean $\theta$-value.
For large $\theta$, the linear term dominates in $\sqrt{\theta^2+1} \approx \theta$, so that
\[
  L_A \approx a\int_{\theta_1}^{\theta_2}\theta\,d\theta
  = \frac{a}{2}(\theta_2^2 - \theta_1^2)
  = \frac{a}{2}(\theta_2+\theta_1)\Delta\theta.
\]
With $\Delta\theta = \frac{2\pi}{3}$ and increasing $\bar\theta = \frac{\theta_1+\theta_2}{2}$,
linear growth in $\bar\theta$ results. The difference of two
consecutive arc lengths is therefore approximately
$a \cdot 6\pi \cdot \frac{2\pi}{3} = a \cdot 4\pi^2 \approx 8.0$ per
index period (from $n$ to $n+6$). Numerically: $25.032 - 17.048 = 7.985$,
$33.025 - 25.032 = 7.993$ etc.\ -- this convergence confirms the asymptotics.
\end{bemerkung}


% ─────────────────────────────────────────────────────────────────────────────
\section{Method B: Arc Length via the $p$-Parametrisation}
\label{sec:methode_B}
% ─────────────────────────────────────────────────────────────────────────────

\subsection{Derivation of the Integral}

In Chapter~\ref{sec:bogenlaenge} the arc length integral in the
$p$-parametrisation was derived. Since $r(p) = \frac{6}{\pi\sqrt{3}}\sqrt{p}$ and
$\theta(p) = \frac{\pi}{\sqrt{3}}\sqrt{p}$, the computation of the
arc length element yields:
\[
\left(\frac{dr}{dp}\right)^2 = \frac{3}{\pi^2 p},
\qquad
r^2\!\left(\frac{d\theta}{dp}\right)^2 = 1,
\]
and hence:
\begin{equation}
\label{eq:bogenlaenge_p}
L_B(p_1, p_2)
= \int_{p_1}^{p_2}
  \frac{1}{\pi}\sqrt{\frac{3}{p} + \pi^2}\,dp.
\end{equation}

The points in the $p$-parametrisation correspond to the function values $p = k(n)$:

\begin{center}
\begin{tabular}{ccccc}
\toprule
$n_1$ & $p_1 = k(n_1)$ & $n_2$ & $p_2 = k(n_2)$ & $\Delta p = p_2 - p_1$ \\
\midrule
 5 & 11  &  7 & 20  & 9  \\
11 & 46  & 13 & 63  & 17 \\
17 & 105 & 19 & 130 & 25 \\
23 & 188 & 25 & 221 & 33 \\
29 & 295 & 31 & 336 & 41 \\
35 & 426 & 37 & 475 & 49 \\
41 & 581 & 43 & 638 & 57 \\
\bottomrule
\end{tabular}
\end{center}

\begin{bemerkung}[Growth of the $\Delta p$ values]
The values $\Delta p = p_2 - p_1$ are precisely the minor steps of the
difference sequence:
\[
\Delta p_m = \Delta_{\mathrm{Minor}}(m) = 8m + 1, \qquad m = 1, 2, 3, \ldots
\]
For $m=1$: $\Delta p = 9$, for $m=2$: $\Delta p = 17$, and so on.
The function values $p$ grow quadratically in $n$, whereas $\Delta p$
grows only linearly. This implies that for large $p$ the quotient
$\Delta p / p$ tends to zero -- the term $\frac{3}{p}$ in the integrand then becomes
negligible compared to $\pi^2$, and the integral approaches the
Euclidean value $\Delta p$.
\end{bemerkung}

\subsection{Numerical Results for Method B}

\begin{table}[H]
\centering
\caption{Arc length Method~B: numerical integral
  $L_B = \int_{p_1}^{p_2} \frac{1}{\pi}\sqrt{\frac{3}{p}+\pi^2}\,dp$.}
\renewcommand{\arraystretch}{1.2}
\begin{tabular}{cccccc}
\toprule
$(n_1, n_2)$ & $p_1 = k(n_1)$ & $p_2 = k(n_2)$ & $\Delta p$ & $L_B$ \\
\midrule
$( 5,  7)$ &  11 &  20 &  9 &  9.09039285 \\
$(11, 13)$ &  46 &  63 & 17 & 17.04772970 \\
$(17, 19)$ & 105 & 130 & 25 & 25.03243825 \\
$(23, 25)$ & 188 & 221 & 33 & 33.02456944 \\
$(29, 31)$ & 295 & 336 & 41 & 41.01977350 \\
$(35, 37)$ & 426 & 475 & 49 & 49.01654429 \\
$(41, 43)$ & 581 & 638 & 57 & 57.01422182 \\
\bottomrule
\end{tabular}
\end{table}


% ─────────────────────────────────────────────────────────────────────────────
\section{Direct Comparison: Method A vs.\ Method B}
\label{sec:methoden_vergleich}
% ─────────────────────────────────────────────────────────────────────────────

\begin{table}[H]
\centering
\caption{Comparison of arc lengths: Method~A (closed-form $\theta$-formula)
  vs.\ Method~B ($p$-integral). The deviation $\delta = (L_B - L_A)/L_A$
  is always positive and decreases monotonically towards zero.}
\renewcommand{\arraystretch}{1.25}
\begin{tabular}{cccccc}
\toprule
$(n_1, n_2)$ & $L_A$ & $L_B$ & $L_B - L_A$ & $\delta$ (\%) \\
\midrule
$( 5,  7)$ &  9.09026487 &  9.09039285 & +0.00012798 & +0.00141 \\
$(11, 13)$ & 17.04771118 & 17.04772970 & +0.00001852 & +0.00011 \\
$(17, 19)$ & 25.03243246 & 25.03243825 & +0.00000579 & +0.00002 \\
$(23, 25)$ & 33.02456693 & 33.02456944 & +0.00000251 & +0.00001 \\
$(29, 31)$ & 41.01977219 & 41.01977350 & +0.00000131 & +0.00000 \\
$(35, 37)$ & 49.01654352 & 49.01654429 & +0.00000077 & +0.00000 \\
$(41, 43)$ & 57.01422134 & 57.01422182 & +0.00000048 & +0.00000 \\
\bottomrule
\end{tabular}
\end{table}

\begin{proposition}[Agreement of the methods]
\label{prop:methoden_uebereinstimmung}
The two computational methods yield numerically almost identical results.
The relative deviation $\delta = (L_B - L_A)/L_A$ is always positive and converges
monotonically to zero:
\[
\delta(5 \to 7) = 0.00141\,\%,\quad
\delta(11 \to 13) = 0.00011\,\%,\quad
\delta(41 \to 43) < 0.000005\,\%.
\]
\end{proposition}

\begin{proof}[Explanation of the small positive deviation]
The deviation arises from the fact that the $p$-parametrisation is only
\emph{asymptotically} arc-length-preserving: the inversion $p = L(\theta)$ holds exactly
for large $\theta$, not for small ones. For the first interval ($p_1 = 11$,
$p_2 = 20$), the correction is still measurable; for $p \gtrsim 100$, it is
negligible for all practical purposes.

Analytically, this can be understood as follows: the arc length element in the
$p$-parametrisation is (as derived in Chapter~\ref{sec:bogenlaenge})
\[
ds = \frac{1}{\pi}\sqrt{\frac{3}{p}+\pi^2}\,dp
   = \sqrt{1 + \frac{3}{\pi^2 p}}\,dp,
\]
which shows that the speed converges to $1$ as $p \to \infty$.
The term $\frac{3}{\pi^2 p}$ provides the positive residual deviation; it is
approximately $0.028$ for $p_1 = 11$, but only approximately $0.0005$ for $p_1 = 581$.
\end{proof}

\subsection{Why the Equidistant Parametrisation Works}

The central conceptual statement of this comparison is that the
integer values $p = k(n)$ of the quadratic sequence already constitute a very good
approximation to the arc length values of the Archimedean spiral. Specifically:
the step $\Delta p = p_2 - p_1 = \Delta_{\mathrm{Minor}}(m)$ is numerically
almost equal to the arc length $L_A$ of the corresponding spiral arc:

\begin{table}[H]
\centering
\caption{Comparison $\Delta p$ (arithmetic step) vs.\  $L_A$ (arc length).
  The ratio $L_A / \Delta p$ converges to~1.}
\renewcommand{\arraystretch}{1.2}
\begin{tabular}{cccccc}
\toprule
$(n_1, n_2)$ & $\Delta p$ & $L_A$ & $L_A / \Delta p$ & Deviation \\
\midrule
$( 5,  7)$ &  9 &  9.09026 & 1.01003 & $+1.00\,\%$ \\
$(11, 13)$ & 17 & 17.04771 & 1.00280 & $+0.28\,\%$ \\
$(17, 19)$ & 25 & 25.03243 & 1.00130 & $+0.13\,\%$ \\
$(23, 25)$ & 33 & 33.02457 & 1.00074 & $+0.07\,\%$ \\
$(29, 31)$ & 41 & 41.01977 & 1.00048 & $+0.05\,\%$ \\
$(35, 37)$ & 49 & 49.01654 & 1.00034 & $+0.03\,\%$ \\
$(41, 43)$ & 57 & 57.01422 & 1.00025 & $+0.02\,\%$ \\
\bottomrule
\end{tabular}
\end{table}

\begin{bemerkung}[Geometric significance]
This table is the numerical centrepiece of the entire construction: it shows
that the arithmetic steps $\Delta p = 8m+1$ of the minor difference sequence
are not only integers, but also increasingly accurate approximations to the
arc lengths of the corresponding spiral arcs. As $m \to \infty$,
$L_A / \Delta p \to 1$ -- the discrete arithmetic structure of the sequence $k(n)$
and the continuous geometry of the Archimedean spiral merge
asymptotically into a single description.

Put differently: the construction succeeds so elegantly because the
function values $k(n)$ are not only integers that enforce the correct spiral constant,
but simultaneously serve as approximate arc length parameters of the
resulting spiral. This is a direct consequence
of the asymptotic arc length formula $L(\theta) \sim \frac{a}{2}\theta^2$ and
the choice $p = k(n) \sim \frac{a}{2}\theta_n^2$.
\end{bemerkung}
